\documentclass[aps,onecolumn,nofootinbib,pra]{article}

\usepackage{../../spec_files/arxiv}
\usepackage{amsmath,amsfonts,amssymb,amsthm,enumerate,times}
\usepackage{mathtools}
\usepackage[usenames,dvipsnames]{color}
\usepackage{hyperref}
\hypersetup{
    breaklinks,
    colorlinks,
    linkcolor=gray,
    citecolor=gray,
    urlcolor=gray,
    pdftitle={},
    pdfauthor={Alex Goessmann}
}

\usepackage{tikz}
\usepackage{graphicx}
\usepackage{float}
\usepackage{comment}
\usepackage{csquotes}
\usepackage{braket}
\usepackage{listings}
\usepackage{verbatim}
\usepackage{etoolbox}
\usepackage[utf8]{inputenc}
\usepackage[english]{babel}
\usepackage[T1]{fontenc}
\usepackage{titlesec}
\usepackage{fancyhdr}
\usepackage{bbm}
\usepackage{bm}
\usepackage{algpseudocode}
\usepackage{algorithm}
\usepackage{lipsum}

\newtheorem{remark}{Remark}
\newtheorem{theorem}{Theorem}
\newtheorem{lemma}{Lemma}
\newtheorem{corollary}{Corollary}
\newtheorem{definition}{Definition}
\newtheorem{example}{Example}

\pretolerance=500
\tolerance=100
\emergencystretch=10pt

% Bibliography
\DeclareUnicodeCharacter{FB01}{fi}
\usepackage[round]{natbib}
\usepackage{wasysym}

\newcommand{\red}[1]{\textcolor{red}{#1}}

\usepackage{import}
\subimport{../../macros/}{all-macros}

\usepackage{minted}

\begin{document}
    \title{Tensor-network based inference for the event calculus}

    \maketitle
    \date{\today}

    \begin{abstract}
        We investigate the tensor-based inference scheme in \cite{tsilionis_tensor-based_2024} exploiting the law of inertia in the event calculus.
        We derive representation of the laws of inertia as tensor network contraction equations exploiting the basis calculus scheme in \cite{goessmann_tensor_2026}.
    \end{abstract}


    \section{Event calculus}

    We introduce two term variables:
    \begin{itemize}
        \item $f$ fluent, enumerated by $\indvariableof{f}$
        \item $t$ time stamps, enumerated by $\indvariableof{t}$
    \end{itemize}

    Further the function:
    \begin{itemize}
        \item $suc(t)$ successor of the time stamp
    \end{itemize}

    And the predicates:
    \begin{itemize}
        \item $h(f,t)$ fluent $f$ holds at time stamp $t$
        \item $e(f,t)$ fluent $f$ ends at time stamp $t$
        \item $s(f,t)$ fluent $f$ starts at time stamp $t$
    \end{itemize}


    \section{Laws of inertia}

    "If a fluent holds or starts at a time $t$, and it is not terminated at time $t$, then it holds also at the successing time stamp $suc(t)$:"
    \begin{align*}
        \big((h(f,t)  \lor s(f,t)) \land \lnot e(f,t) \big)
        \rightarrow h(f,suc(t))
    \end{align*}

    Exploiting basis calculus, these pose the contraction equation on the basis encodings of the predicates and functions (which specify the FOL worlds):
    \stantikzpic{

        \drawtensorblock{9}{8}{10}{$\impbincon$}

        \drawtdwire{0}{7}{$\headvariableof{\lnot e}$}
        \drawsmalltensor{0}{4}{$\bencodingof{\lnot}$}
        \drawtdwire{0}{3}{$\headvariableof{e}$}
        \drawtensorblock{0}{0}{2}{$\bencodingof{e}$}

        \drawtdwire{7.5}{7}{$\headvariableof{\lor}$}
        \drawtensorblock{7.5}{4}{3}{$\bencodingof{\lor}$}

        \drawtdwire{5}{3}{$\headvariableof{s}$}
        \drawtensorblock{5}{0}{2}{$\bencodingof{s}$}

        \drawtdwire{10}{3}{$\headvariableof{h}$}
        \drawtensorblock{10}{0}{2}{$\bencodingof{h}$}

        \drawtdwire{18}{7}{$\headvariableof{h}$}
        \drawtensorblock{18}{4}{2}{$\bencodingof{h}$}
        \drawtdwire{19.5}{3}{$\indvariableof{suc(t)}$}
        \drawtensorblock{19.5}{0}{2}{$\bencodingof{suc}$}

        \draw[->-] (5,-4) to[bend right=-20] (-1.5,-1);
        \draw[->-] (5,-4) to[bend right=-5] (3.5,-1);
        \draw[->-] (5,-4) to[bend right=15] (8.5,-1);
        \draw[->-] (5,-4) to[bend right=25] (16.5,3);
        \drawvariabledot{5}{-4} % t
        \drawtdwire{5}{-4}{$\indvariableof{f}$}

        \draw[->-] (10,-4) to[bend right=-20] (1.5,-1);
        \draw[->-] (10,-4) to[bend right=-5] (6.5,-1);
        \draw[->-] (10,-4) to[bend right=15] (11.5,-1);
        \draw[->-] (10,-4) to[bend right=18] (19.5,-1);
        \drawvariabledot{10}{-4} % t
        \drawtdwire{10}{-4}{$\indvariableof{t}$}


        \drawtextnode{25}{0}{$=$}
        \drawtensorblock{30}{0}{2}{$\ones$}
        \drawtdwire{28.5}{-1}{$\indvariableof{f}$}
        \drawtdwire{31.5}{-1}{$\indvariableof{t}$}
    }

    When choosing the time order on the time stamps, the basis encoding of $suc$ is (called $U$ in \cite{tsilionis_tensor-based_2024})
    \begin{align*}
        \bencodingofat{suc}{\indexedindvariableof{suc(t)},\indexedindvariableof{t}}
        = \begin{cases}
              1 & \ifspace \indindexof{suc(t)} = \indindexof{t}+1 \\
              0 & \elsetext \, .
        \end{cases}
    \end{align*}


    \section{Inference}

    Let us consider that the groundings to the predicates $s[\indvariableof{f},\indvariableof{t}]$ and $t[\indvariableof{f},\indvariableof{t}]$ and to the function $suc$ are known, and consider infering the least tensor $h[\indvariableof{f},\indvariableof{t}]$ respecting the laws of inertia:
    \begin{align*}
        \argmin_{h[\indvariableof{f},\indvariableof{t}]} \contraction{h[\indvariableof{f},\indvariableof{t}]}
        \quad \text{subject to the laws of inertia holding for given}\quad s[\indvariableof{f},\indvariableof{t}],t[\indvariableof{f},\indvariableof{t}]
    \end{align*}

    To solve the contraction equation we contract with any $\onehotmapof{\indindexof{t}}$ for $\indindexofin{t}$ to get the equation:

    \stantikzpic{

        \drawtensorblock{9}{8}{10}{$\impbincon$}

        \drawtdwire{0}{7}{$\headvariableof{\lnot e}$}
        \drawsmalltensor{0}{4}{$\bencodingof{\lnot}$}
        \drawtdwire{0}{3}{$\headvariableof{e}$}
        \drawtensorblock{0}{0}{2}{$\bencodingof{e}$}

        \drawtdwire{7.5}{7}{$\headvariableof{\lor}$}
        \drawtensorblock{7.5}{4}{3}{$\bencodingof{\lor}$}

        \drawtdwire{5}{3}{$\headvariableof{s}$}
        \drawtensorblock{5}{0}{2}{$\bencodingof{s}$}

        \drawtdwire{10}{3}{$\headvariableof{h}$}
        \drawtensorblock{10}{0}{2}{$\bencodingof{h}$}

        \drawtdwire{18}{7}{$\headvariableof{h}$}
        \drawtensorblock{18}{4}{2}{$\bencodingof{h}$}
        \drawtdwire{19.5}{3}{$\indvariableof{suc(t)}$}
        \drawtensorblock{19.5}{0}{1.5}{$\onehotmapof{\indindexof{t}+1}$}

        \draw[->-] (5,-4) to[bend right=-20] (-1.5,-1);
        \draw[->-] (5,-4) to[bend right=-5] (3.5,-1);
        \draw[->-] (5,-4) to[bend right=15] (8.5,-1);
        \draw[->-] (5,-4) to[bend right=25] (16.5,3);
        \drawvariabledot{5}{-4} % t
        \drawtdwire{5}{-4}{$\indvariableof{f}$}

        \draw[->-] (1.5,-6) -- (1.5,-1);
        \drawsmalltensor{1.5}{-7}{$\onehotmapof{\indindexof{t}}$}
        \draw[->-] (6.5,-6) -- (6.5,-1);
        \drawsmalltensor{6.5}{-7}{$\onehotmapof{\indindexof{t}}$}
        \draw[->-] (11.5,-6) -- (11.5,-1);
        \drawsmalltensor{11.5}{-7}{$\onehotmapof{\indindexof{t}}$}

        \drawtextnode{25}{0}{$=$}
        \drawtensorblock{30}{0}{1}{$\ones$}
        \drawtdwire{30}{-1}{$\indvariableof{t}$}
    }

    We solve them iteratively by computing the vectors on the right hand side and write them into the slices of $h$ as:
    \stantikzpic{

        \drawtensorblock{3.75}{8}{5}{$\land$}

        \drawtdwire{0}{7}{$\headvariableof{\lnot e}$}
        \drawsmalltensor{0}{4}{$\bencodingof{\lnot}$}
        \drawtdwire{0}{3}{$\headvariableof{e}$}
        \drawtensorblock{0}{0}{2}{$\bencodingof{e}$}

        \drawtdwire{7.5}{7}{$\headvariableof{\lor}$}
        \drawtensorblock{7.5}{4}{3}{$\bencodingof{\lor}$}

        \drawtdwire{5}{3}{$\headvariableof{s}$}
        \drawtensorblock{5}{0}{2}{$\bencodingof{s}$}

        \drawtdwire{10}{3}{$\headvariableof{h}$}
        \drawtensorblock{10}{0}{2}{$\bencodingof{h}$}



        \draw[->-] (5,-4) to[bend right=-20] (-1.5,-1);
        \draw[->-] (5,-4) to[bend right=-5] (3.5,-1);
        \draw[->-] (5,-4) to[bend right=15] (8.5,-1);

        %\draw[->-] (5,-4) to[bend right=25] (16.5,3);
        \drawvariabledot{5}{-4} % t
        \drawtdwire{5}{-4}{$\indvariableof{f}$}

        \draw[->-] (1.5,-6) -- (1.5,-1);
        \drawsmalltensor{1.5}{-7}{$\onehotmapof{\indindexof{t}}$}
        \draw[->-] (6.5,-6) -- (6.5,-1);
        \drawsmalltensor{6.5}{-7}{$\onehotmapof{\indindexof{t}}$}
        \draw[->-] (11.5,-6) -- (11.5,-1);
        \drawsmalltensor{11.5}{-7}{$\onehotmapof{\indindexof{t}}$}



        \drawtextnode{-8}{0}{$\algdefsymbol$}
        \drawtensorblock{-15}{0}{3}{${h}$}
        \drawtdwire{-13.5}{-1}{$\indvariableof{t}$}
        \drawtdwire{-16.5}{-1}{$\indvariableof{f}$}
        \drawtensorblock{-13.5}{-4}{1.5}{$\onehotmapof{\indindexof{t}+1}$}
    }

    Note that we exploit in the solution the consecutive dependence structure of the equations, induced by the $\bencodingof{suc}$ matrix.

    \section{Generic solution strategies}

    The solution of contraction equations can in general be done by simplifying them:
    \begin{itemize}
        \item Close some open legs with arbitrary tensors: When choosing basis tensors (or more general normed tensors), the right hand sight is still a $\ones$ tensor.
        We used this trick to derive sequences of simpler equations for each time stamp.
        \item Contract only a subset of the left hand side: We could restrict by removing $\bencodingof{s}$ and $\bencodingof{h}$ to the two original laws of inertia in \cite{tsilionis_tensor-based_2024}.
    \end{itemize}

    \bibliographystyle{plainnat}
    \bibliography{../../references.bib}


\end{document}